O monitoramento de infraestruturas críticas, como o Estreito de Ormuz, constitui um desafio complexo no campo dos sistemas distribuídos. O estreito é uma das rotas marítimas mais estratégicas do mundo, responsável por aproximadamente 20\% do tráfego global de petróleo \cite{tanenbaum2021redes}. A necessidade de rastrear continuamente embarcações, detectar anomalias ambientais e despachar recursos de resposta a emergências exige uma arquitetura de software robusta, de alta disponibilidade e capaz de operar mesmo diante de falhas parciais na infraestrutura.

Sistemas baseados em modelo de \textit{broker} centralizado, como o FarmNode desenvolvido no projeto anterior, introduzem um ponto único de falha (\textit{Single Point of Failure} -- SPOF) crítico: se o servidor central torna-se inoperante, toda a infraestrutura de monitoramento e salvamento marítimo é interrompida imediatamente. Adicionalmente, a alocação concorrente de recursos escassos --- como drones de patrulhamento disputados por múltiplas ocorrências simultâneas --- pode resultar em inanição de tarefas ou impasse (\textit{deadlock}) distribuído.

Para mitigar esses problemas, este trabalho apresenta o \textbf{HormuzNet}, um sistema de monitoramento marítimo descentralizado baseado em uma malha cooperativa de \textit{brokers}. O HormuzNet elimina SPOFs ao distribuir o processamento por meio de uma arquitetura híbrida que combina o modelo cliente-servidor em dois níveis com uma malha ponto a ponto (P2P). A consistência de dados e a ordenação global de eventos são garantidas pelos Relógios Lógicos de Lamport \cite{lamport1978time}, enquanto a exclusão mútua distribuída resolve a concorrência pelo despacho de drones.

O sistema foi implementado em Go com protocolos nativos de rede: UDP/IP Multicast para telemetria redundante, TCP/IP persistente com lógica de \textit{fallback} para coordenação de drones autônomos, e WebSocket para visualização tática em tempo real. A infraestrutura foi validada com 34 contêineres Docker distribuídos entre duas máquinas físicas interconectadas por rede LAN Gigabit.

Os resultados demonstram que a malha de \textit{brokers} reage autonomamente a falhas por meio de \textit{Ring Failover} e eleição Bully, realizando o despacho de drones com latência média inferior a 1,5\,s e sem inconsistências nas filas de prioridade distribuídas.
